\documentclass[11pt]{article}
% =========================================================
%  學生講義 共用前言（以 XeLaTeX 編譯）
%  需要套件：xeCJK, tcolorbox（其餘多在 BasicTeX 內）
% =========================================================
\usepackage[a4paper,margin=2.3cm]{geometry}
\usepackage{amsmath,amssymb}
\usepackage{xcolor}
\usepackage{graphicx}

% ---- 中文字型（macOS 系統字型）----
\usepackage{xeCJK}
\setCJKmainfont{Songti TC}      % 內文：宋體
\setCJKsansfont{Heiti TC}       % 標題/強調：黑體
\xeCJKsetup{AutoFakeBold=2.2, AutoFakeSlant=0.2}

% ---- 版面 ----
\setlength{\parindent}{0pt}
\setlength{\parskip}{0.55em}
\linespread{1.15}
\renewcommand{\baselinestretch}{1.15}

% ---- 顏色 ----
\definecolor{cBlue}{HTML}{1f6feb}
\definecolor{cGray}{HTML}{6b7280}
\definecolor{cGrayBg}{HTML}{f3f4f6}
\definecolor{cPurp}{HTML}{7a3ff2}
\definecolor{cPurpBg}{HTML}{f5f1ff}

% ---- 標註框 ----
\usepackage[most]{tcolorbox}

% 就地補充新概念（灰底）：\begin{補充框}{標題}
\newtcolorbox{補充框}[1]{
  enhanced, breakable, arc=2pt, boxrule=0.6pt,
  colback=cGrayBg, colframe=cGray,
  left=9pt,right=9pt,top=7pt,bottom=7pt,
  before upper={\textbf{\sffamily #1}\par\vspace{3pt}},
}

% 延伸 / 開放問題（紫色左邊條）：\begin{延伸框}{標題}
\newtcolorbox{延伸框}[1]{
  enhanced, breakable, arc=2pt, boxrule=0pt,
  colback=cPurpBg, colframe=cPurpBg,
  borderline west={3pt}{0pt}{cPurp},
  left=10pt,right=9pt,top=7pt,bottom=7pt,
  before upper={\textbf{\sffamily 延伸閱讀｜#1}\par\vspace{3pt}},
}

% ---- 圖：置中、底下小字說明 ----
% \fig{檔名}{寬度}{說明}
\newcommand{\fig}[3]{%
  \begin{center}
  \includegraphics[width=#2\linewidth]{#1}\\[2pt]
  {\small\color{cGray} #3}
  \end{center}}

% ---- 章節標題用黑體 ----
\newcommand{\h}[1]{\sffamily #1}


% 讓羅馬數字 Ⅰ（U+2160）由中文字型排版，避免落到無此字符的拉丁字型
\xeCJKDeclareCharClass{CJK}{`Ⅰ}

\begin{document}

\begin{center}
{\sffamily\Huge\bfseries 選物I-6 萬有引力}\\[3pt]
{\sffamily\large 普通高中 選修物理Ⅰ（力學一）・學生講義}
\end{center}
\vspace{0.6em}

必修課程已定性提到：牛頓把「天上」與「地上」的力學統一在同一套規則之下——讓蘋果落地的力，與讓月亮繞地球運轉的力，是同一種力。本章把這個結論寫成一條精確的定律：宇宙中任意兩個有質量的物體，都以 $F=\dfrac{GMm}{r^2}$ 互相吸引。這稱為\textbf{萬有引力定律（law of universal gravitation）}。

由這一條定律，可以：算出地表的重力加速度 $g$、說明衛星如何繞地球運行而不落下、並把克卜勒第三定律 $T^2\propto r^3$ 由力學推導出來。本章是必修「天上地下同一套力學」的\textbf{定量加強版}，也是後續學習重力位能、力學能守恆與角動量的基礎。

\medskip
本章主線：

\begin{center}
萬有引力定律 $F=\dfrac{GMm}{r^2}$ $\rightarrow$ 地表重力 $g=\dfrac{GM}{R^2}$ $\rightarrow$ 衛星圓軌道（萬有引力＝向心力）$\rightarrow$ 由圓軌道推導克卜勒第三定律 $T^2\propto r^3$
\end{center}

\medskip
\textbf{貫穿全章的提醒：}式子裡的 $r$ 是\textbf{兩物體質心之間的距離}。算地表物體所受地球引力時，$r$ 要用\textbf{地球半徑}（從地心算起），不是 0，也不是離地面的高度。把 $r$ 取錯是本章最常見的錯誤。

\section*{\sffamily 6-1\quad 萬有引力定律}

\subsection*{\sffamily A.\ 從克卜勒定律到一條普適的力}

牛頓的突破不在於「發現重力」——物體會落下，這是日常可見的。他的突破在於提出：\textbf{讓蘋果落地的力，與讓月亮繞地球運轉的力，是同一種力}。在此之前，「天上」（被視為完美、永恆的圓周運動）與「地上」（會落下、會變化的物體）被當作服從不同規則的兩個領域。牛頓將兩者合而為一：天上與地上，受同一條引力定律支配。

牛頓進一步推測這個力的形式。由克卜勒第三定律（必修已學：$T^2\propto a^3$）配合圓周運動的向心力公式反推，可得到這個力必須與\textbf{距離平方成反比}（這正是 6-4 要進行的推導，此處先承認結論）。由此寫下萬有引力定律：

\begin{補充框}{萬有引力定律 $F=\dfrac{GMm}{r^2}$}
宇宙中任意兩個質點，彼此以\textbf{沿兩者連線方向的吸引力}互相作用，大小為
$$\boxed{\;F=\dfrac{GMm}{r^{2}}\;}$$
其中 $M$、$m$ 為兩質點質量，$r$ 為兩者\textbf{質心間的距離}，$G$ 為\textbf{萬有引力常數（gravitational constant）}：
$$G\approx 6.67\times10^{-11}\ \text{N·m}^2/\text{kg}^2.$$
這個力是一對\textbf{作用力與反作用力}：$M$ 拉 $m$，$m$ 也以等大、反向的力拉 $M$。
\end{補充框}

\fig{si6-law}{0.95}{圖 6-1：左——兩質量間的萬有引力沿連線、互相吸引、等大反向。右——平方反比關係：距離變 2 倍，引力變為 $\tfrac14$；距離變 3 倍，引力變為 $\tfrac19$。}

讀懂這條式子，有三個要點：

\begin{enumerate}\setlength{\itemsep}{2pt}
\item \textbf{平方反比（inverse-square）}：距離變 2 倍，引力變為 $\dfrac{1}{2^2}=\dfrac14$；距離變 3 倍，引力變為 $\dfrac19$。引力衰減得快，但\textbf{不會真正變成零}——這是「萬有」二字的由來。
\item \textbf{與兩質量的乘積成正比}：兩個質量都加倍，引力變 4 倍。
\item \textbf{$G$ 極小}：因此日常物體之間的引力小到無法察覺；要產生明顯的引力，至少要有一個如行星般巨大的質量。例如兩個 60 kg 的人相距 1 m，彼此引力僅約 $2.4\times10^{-7}$ N，比體重小約 25 億倍。
\end{enumerate}

\begin{補充框}{$r$ 取哪裡？可以把大球當成一個點嗎？}
\textbf{注意：}$r$ \textbf{不是}兩物體表面之間的距離，而是\textbf{質心到質心}的距離。計算地表物體所受地球引力時，$r$ 要用\textbf{地球半徑}（從地心算起），不是 0。

那麼，把地球這種大球當成一個點來算，合理嗎？牛頓證明過：一個質量\textbf{均勻的球體（或球殼）}，對外部物體的引力，\textbf{恰好等於把它的全部質量集中在球心}所產生的引力。因此把地球視為「質量集中於地心的一個質點」，從地心起算距離 $R$，是\textbf{嚴格正確}的，並非近似。
\end{補充框}

\begin{補充框}{$G$ 與 $g$ 是兩個不同的量}
$G$（萬有引力常數）與 $g$（地表重力加速度）容易混淆，兩者其實不同：
\begin{itemize}\setlength{\itemsep}{1pt}
\item \textbf{$G$} 是\textbf{全宇宙通用}的基本常數，數值 $6.67\times10^{-11}$，描述「引力有多強」，到哪裡都一樣。
\item \textbf{$g$} 是\textbf{地球專屬}的重力加速度，數值約 $9.8\ \text{m/s}^2$，是萬有引力在地表的具體後果。6-2 將證明 $g$ 是由 $G$ 算出來的。
\end{itemize}
$G$ 的數值由卡文迪西（Cavendish）以扭秤測出。在此之前，人類無法分離出 $G$ 與地球質量 $M$，只能得到兩者綁在一起的組合。
\end{補充框}

\section*{\sffamily 6-2\quad 地表重力與重力加速度}

\subsection*{\sffamily A.\ 把 $W=mg$ 與 $F=\dfrac{GMm}{r^2}$ 接起來}

一個質量 $m$ 的物體放在地表，受到地球（質量 $M$、半徑 $R$）的萬有引力，這就是平常所說的「重力 $W$」。物體與地心的距離是地球半徑 $R$，所以
$$W=\dfrac{GMm}{R^{2}}.$$
而我們也習慣寫 $W=mg$。兩式描述的是\textbf{同一個力}，令其相等：
$$mg=\dfrac{GMm}{R^{2}}\;\Longrightarrow\;\boxed{\,g=\dfrac{GM}{R^{2}}\,}$$

兩邊的 $m$ 消去了。這說明\textbf{重力加速度 $g$ 與物體本身的質量無關}——在真空中，羽毛與鉛球落下的快慢相同。$g$ 完全由地球的 $M$ 與 $R$ 決定。把地球的 $M\approx5.97\times10^{24}\ \text{kg}$、$R\approx6.37\times10^{6}\ \text{m}$ 與 $G$ 代入，算得的正是 $g\approx9.8\ \text{m/s}^2$。可見 $g$ 不是憑空給定的數值，而是可由 $G$、$M$、$R$ 算出來的。

\begin{補充框}{$g=\dfrac{GM}{R^2}$ 的意義，以及如何「秤地球」}
$G$ 與 $g$ 雖是不同的量，卻\textbf{被同一個力連在一起}：地表的重力既可寫成 $mg$，也可寫成 $\dfrac{GMm}{R^2}$，因為兩者是同一個力。把兩種寫法畫上等號，再消去物體質量 $m$，就「算」出 $g=\dfrac{GM}{R^2}$。

可類比為：$G$ 像「全國統一的稅率公式」（人人共用），$g$ 像「某一間店在特定地點、特定規模下實際要繳的稅額」（因店而異）。

反過來，把 $g=\dfrac{GM}{R^2}$ 解出 $M$：
$$M=\dfrac{gR^{2}}{G}.$$
$g$、$R$ 都可在地表量得，$G$ 由卡文迪西實驗測得，代入即得\textbf{整顆地球的質量} $M\approx5.97\times10^{24}\ \text{kg}$。人類未曾離開地面，便算出了腳下星球的質量。
\end{補充框}

\subsection*{\sffamily B.\ $g$ 會隨高度與緯度改變}

\textbf{隨高度變化}：在離地面高度 $h$ 處，與地心的距離變成 $R+h$，所以
$$g_h=\dfrac{GM}{(R+h)^{2}}.$$
離地越高，$g$ 越小，但不會降到零。注意分母是 $(R+h)^2$，不是 $h^2$：飛機飛在約 10 km 高空，相較於 $R\approx6370$ km，$g$ 幾乎不變；要使 $g$ 明顯變小，需上升到「以地球半徑為單位」的高度。例如要使 $g$ 剩地表的 $\tfrac14$，需上升一個地球半徑（約 6370 km），此時 $R+h=2R$。

\fig{si6-surface-g}{0.66}{圖 6-2：地表重力加速度隨高度遞減，$g=\dfrac{GM}{(R+h)^2}$。地表（$r=R$）$g\approx9.8$；上升約一個地球半徑（$r=2R$）時 $g$ 剩約 $\tfrac14$。曲線在地表附近變化緩慢。}

\textbf{隨緯度變化}：實測 $g$ 在赤道約 $9.78$，在兩極約 $9.83$。原因有二：(1) 地球並非完美的球，\textbf{赤道略鼓}，赤道地表離地心較遠，$g$ 較小；(2) 地球自轉，赤道上的物體需要一部分向心力，由重力提供，使站在磅秤上量到的\textbf{視重略小}。

\begin{補充框}{質量（mass）與重量（weight）的區別}
\begin{itemize}\setlength{\itemsep}{1pt}
\item \textbf{質量} 是物體所含物質的多寡、抵抗加速難易程度的度量，\textbf{到哪裡都一樣}：一塊 1 kg 的鐵帶到月球，仍是 1 kg。
\item \textbf{重量} 是物體所受\textbf{重力的大小} $W=mg$，\textbf{會隨所在地的 $g$ 改變}：月球 $g$ 約為地球的 $\tfrac16$，同一塊鐵在月球的重量只剩約 $\tfrac16$，但質量不變。
\end{itemize}
\textbf{注意：}磅秤量到的其實是重量（更精確地說是支撐物體的正向力），卻以「公斤」標示，這是生活上的便宜用法，不應因此把重量與質量混為一談。
\end{補充框}

\section*{\sffamily 6-3\quad 行星與人造衛星的運動}

\subsection*{\sffamily A.\ 核心觀念：萬有引力就是向心力}

衛星為何能繞地球運行而不落下？關鍵在於一句話：\textbf{讓衛星做圓周運動所需的向心力，正是由地球對它的萬有引力提供}。

在真空中（忽略大氣），衛星只受萬有引力一個力。它因慣性而傾向沿直線飛離，萬有引力則不斷把它的路徑「拉彎」，使其彎成一個圓。設圓軌道半徑為 $r$、衛星速率為 $v$，將上述關係寫成式子：
$$\underbrace{\dfrac{GMm}{r^{2}}}_{\text{萬有引力}}=\underbrace{\dfrac{mv^{2}}{r}}_{\text{向心力}}.$$
消去衛星質量 $m$（與衛星本身質量無關），解出 $v$：
$$\dfrac{GMm}{r^{2}}=\dfrac{mv^{2}}{r}\;\Longrightarrow\;v^{2}=\dfrac{GM}{r}\;\Longrightarrow\;\boxed{\,v=\sqrt{\dfrac{GM}{r}}\,}$$

\fig{si6-orbit}{0.50}{圖 6-3：衛星圓軌道，萬有引力提供向心力 $\dfrac{GMm}{r^2}=\dfrac{mv^2}{r}$。軌道半徑越大，速率越小：近軌道（藍）較快，遠軌道（綠）較慢。}

由 $v=\sqrt{\dfrac{GM}{r}}$ 可知：\textbf{軌道半徑 $r$ 越大，速率 $v$ 越小}——外圈的衛星運行得較慢。$v$ 不含衛星質量 $m$，因此在同一軌道上，一顆螺絲釘與一座太空站的速率相同。

\textbf{注意：}有一種直覺認為「軌道大、要繞的圈長，所以跑得快」，這與式子相反。$v=\sqrt{GM/r}$ 表示 $r$ 越大，$v$ 越\textbf{小}。低軌道衛星（如國際太空站，約 400 km 高）速率約 $7.7\ \text{km/s}$；高軌道的同步衛星則慢得多。離得越遠，引力越弱，只能、也只需維持較慢的速率。

\begin{補充框}{常用技巧：恆等式 $GM=gR^2$}
許多衛星問題不直接給 $G$、$M$，只給 $g$ 與 $R$。由 $g=\dfrac{GM}{R^2}$ 移項，得
$$GM=gR^2.$$
利用這個恆等式，可把含 $GM$ 的式子（如 $v=\sqrt{GM/r}$）改寫成只含 $g$、$R$、$r$ 的形式，避免代入 $10^{24}$、$10^{-11}$ 等大小懸殊的數字。例如貼近地表的圓軌道（$r\approx R$）速率為
$$v=\sqrt{\dfrac{GM}{R}}=\sqrt{\dfrac{gR^2}{R}}=\sqrt{gR}.$$
\end{補充框}

\subsection*{\sffamily B.\ 第一宇宙速度與同步衛星}

由上式，貼近地表的圓軌道速率為
$$v_1=\sqrt{gR}=\sqrt{\dfrac{GM}{R}}\approx7.9\ \text{km/s}.$$
這稱為\textbf{第一宇宙速度（first cosmic speed）}：物體能「貼著地表繞地球運行」所需的最小速率，也是任何近地軌道衛星的速率量級。低於此速率，物體會落回地面；達到此速率，便能持續環繞。對應的繞行週期約 85 分鐘，與真實低軌衛星（約 90 分鐘繞地球一圈）接近。

\textbf{地球同步衛星（geostationary satellite）}的週期恰等於地球自轉週期（$T=24\ \text{小時}=86400\ \text{s}$）。它位於赤道上空，與地面同步轉動，從地面看彷彿靜止懸停在天空同一點，氣象與通訊衛星多採此軌道。將「萬有引力＝向心力」配合 $v=\dfrac{2\pi r}{T}$，可解出其軌道半徑
$$\dfrac{GM}{r^2}=\dfrac{4\pi^2 r}{T^2}\;\Longrightarrow\;r^3=\dfrac{GM\,T^2}{4\pi^2}\;\Longrightarrow\;r\approx4.2\times10^{7}\ \text{m},$$
離地面約 $3.6\times10^{7}$ m（約 36000 km），與實際同步衛星軌道相符。

\subsection*{\sffamily C.\ 「失重」不是「沒有重力」}

太空人在太空站內漂浮，常被說成「太空中沒有重力」。這是不正確的。在 400 km 高處，$g$ 仍約為地表的 $89\%$，重力幾乎與地面相當。漂浮的原因不是重力消失，而是：\textbf{太空人與太空站以相同的加速度（重力提供的向心加速度）一起自由下落}。兩者同步下落，人對地板（太空站）便沒有壓力，磅秤讀數為零。這稱為\textbf{失重（weightlessness）}，更精確的說法是\textbf{表觀失重}。它與「電梯纜繩斷裂時，人在電梯內會飄起來」是同一回事。

\begin{補充框}{衛星為何不落下？失重的成因}
\textbf{衛星其實一直在「往地球掉」}，只是水平速度夠快，下落的弧線剛好與地球表面一樣彎，於是永遠掉不到地面。可由牛頓設想的思想實驗（牛頓大砲）理解：在高山頂水平發射砲彈，初速越大射程越遠；當初速達到約 $7.9\ \text{km/s}$，砲彈下墜的弧線彎曲程度恰與地表曲率一致，便不再落地，而是環繞地球，成為衛星。衛星與拋出的石頭本質相同，\textbf{都在自由下落}，差別只在水平速度的大小。

由公式看：衛星的向心加速度 $a=\dfrac{v^2}{r}$ 完全由該處的重力 $g_r=\dfrac{GM}{r^2}$ 提供，即 $a=g_r$。也就是說，衛星確實在以該處的 $g_r$ 往地心方向加速（下落），只是同時具有龐大的切線速度，使下落的軌跡閉合成一個圓。

「重量」其實是地板（或磅秤）對人施加的正向力 $N$。站在地面時，地板撐住人，$N=mg$，磅秤讀出體重。當人與地板\textbf{以相同的加速度一起自由下落}時，地板不再撐住人，$N=0$，磅秤讀數歸零，人便飄起來。\textbf{注意：}失重指的是「沒有支撐力、視重為零」，並非「沒有重力」；太空人仍受著巨大的萬有引力，那正是他繞地球運行的原因。
\end{補充框}

\section*{\sffamily 6-4\quad 由萬有引力推導克卜勒第三定律}

必修課程中，克卜勒三大定律是克卜勒由第谷的觀測資料\textbf{歸納}出的經驗規律——他知道「規律是這樣」，但不知道「為什麼」。牛頓的萬有引力定律提供了「為什麼」：只要假設行星受太陽的平方反比引力，三大定律都可由力學\textbf{演繹}出來。以下用最簡單的\textbf{圓軌道}為例，把第三定律 $T^2\propto r^3$ 推導出來。

\subsection*{\sffamily A.\ 推導（圓軌道）}

設行星質量 $m$，繞太陽（質量 $M$）做半徑 $r$、週期 $T$ 的圓周運動，萬有引力提供向心力。

\textbf{第 1 步：寫下「萬有引力＝向心力」。}
$$\dfrac{GMm}{r^{2}}=\dfrac{mv^{2}}{r}.$$

\textbf{第 2 步：把速率 $v$ 換成週期 $T$。}繞一圈的路程是圓周長 $2\pi r$，所花時間是一個週期 $T$，所以
$$v=\dfrac{2\pi r}{T}\;\Longrightarrow\;v^2=\dfrac{4\pi^2 r^2}{T^2}.$$

\textbf{第 3 步：代回。}
$$\dfrac{GMm}{r^{2}}=\dfrac{m}{r}\cdot\dfrac{4\pi^2 r^2}{T^2}=\dfrac{4\pi^2 m r}{T^2}.$$

\textbf{第 4 步：消去 $m$、整理。}兩邊消去行星質量 $m$（故結論與行星本身質量無關），再把 $T^2$、$r^3$ 各歸一邊：
$$\dfrac{GM}{r^{2}}=\dfrac{4\pi^2 r}{T^2}\;\Longrightarrow\;GM\,T^2=4\pi^2 r^3\;\Longrightarrow\;\boxed{\,T^{2}=\dfrac{4\pi^2}{GM}\,r^{3}\,}.$$

\fig{si6-kepler}{0.95}{圖 6-4：左——由「圓軌道引力＝向心力」推得 $T^2=\dfrac{4\pi^2}{GM}r^3$。右——太陽系各行星的 $T^2$ 對 $r^3$ 作圖為一條過原點的直線，斜率為 $k=\dfrac{4\pi^2}{GM}$，即 $T^2\propto r^3$。}

\subsection*{\sffamily B.\ 讀出克卜勒第三定律}

$\dfrac{4\pi^2}{GM}$ 是一個\textbf{常數}（只與中央天體質量 $M$ 有關，與行星無關）。所以
$$T^2=k\,r^3\;\Longrightarrow\;\boxed{\,T^2\propto r^3\,},\qquad k=\dfrac{4\pi^2}{GM}.$$
這正是\textbf{克卜勒第三定律}，而且連比例常數 $k$ 都一併得到。對繞同一個太陽的所有行星，這個 $k$ 都相同，這說明了「離太陽越遠的行星，週期越長」。

\begin{補充框}{讀這條式子要注意兩點}
\begin{itemize}\setlength{\itemsep}{1pt}
\item \textbf{$T^2\propto r^3$ 並非 $T$ 與 $r$ 成正比。}是 $T$ 的平方正比於 $r$ 的立方。$r$ 變 4 倍，$T$ 變 $4^{3/2}=8$ 倍，不是 4 倍。
\item \textbf{比例常數 $k$ 與行星質量無關。}第 4 步已消去行星質量 $m$，$k=\dfrac{4\pi^2}{GM}$ 只與中央天體 $M$ 有關。木星與地球繞太陽，用的是同一個 $k$。
\end{itemize}
此式由圓軌道推得；嚴格的克卜勒第三定律對橢圓軌道也成立，只要把 $r$ 換成橢圓的\textbf{半長軸（semi-major axis）}$a$：$T^2=\dfrac{4\pi^2}{GM}a^3$。圓是離心率為 0 的特殊橢圓，此時 $a$ 即等於半徑 $r$，故圓的結論是一般式的特例。橢圓情形的嚴格證明需用到微積分（大學內容），但高中以圓軌道得到的結論，形式與係數完全正確。
\end{補充框}

\begin{延伸框}{行星軌道為什麼是橢圓，而不是正圓}
本章為了數學簡單，全用圓軌道計算；但克卜勒第一定律指出，真實軌道是\textbf{橢圓}，太陽位於一個焦點上（\textbf{注意：}是焦點，不是中心；橢圓有兩個焦點，太陽佔其一，另一個是空的）。為什麼是橢圓？

\textbf{圓所需的條件非常苛刻。}圓是橢圓的一個特例（離心率 $e=0$）。要維持正圓，行星在每個位置的速率都必須恰到好處——快一點軌道就被甩成往外的橢圓，慢一點就被拉成往內的橢圓。這種精確條件在真實系統中幾乎不會被滿足，因此行星軌道\textbf{一般是橢圓}。地球軌道很接近圓（離心率僅約 0.017），但仍是橢圓。

\textbf{軌道形狀由總力學能 $E$ 決定。}物體在平方反比引力場中，所有可能的軌道恰好是\textbf{圓錐曲線（conic sections）}——以平面切圓錐所得的曲線族：
\begin{itemize}\setlength{\itemsep}{1pt}
\item $E<0$（束縛態）：\textbf{橢圓}（圓是 $E$ 最低的特例）。行星、衛星被綁住反覆繞行。
\item $E=0$（臨界）：\textbf{拋物線}。恰以脫離速度離去，一去不返。
\item $E>0$（非束縛）：\textbf{雙曲線}。速度超過脫離速度，飛掠而過，如某些彗星。
\end{itemize}
行星被太陽束縛（$E<0$），所以走橢圓。橢圓並非被特別挑出，它只是「被綁住」這個最常見情形所對應的形狀。

可補充一點：唯有平方反比力（與正比於距離的彈簧力）會給出\textbf{封閉、不進動}的橢圓軌道（此為伯特蘭定理 Bertrand's theorem）。若力的形式略偏離平方反比，每繞一圈，橢圓長軸會緩慢轉動。因此「行星走固定不動的橢圓、太陽在焦點」這件事，正是重力恰為平方反比的一個敏感印證。

\textbf{注意：}橢圓軌道並非四季的成因。地球軌道幾乎是圓，日地距離一年只差約 3\%，不足以造成四季；四季的主因是\textbf{地軸傾斜}。
\end{延伸框}

\section*{\sffamily 6-5\quad 本章重點整理}

\begin{補充框}{一頁複習（公式與關鍵結論）}
\begin{itemize}\setlength{\itemsep}{2pt}
\item \textbf{萬有引力定律}：$F=\dfrac{GMm}{r^2}$，沿連線、互相吸引、成對出現，$G\approx6.67\times10^{-11}$。$r$ 是\textbf{質心間距離}；均勻球可視為質量集中於球心。平方反比：$r$ 變 2 倍，$F$ 變 $\tfrac14$。
\item \textbf{地表重力}：$g=\dfrac{GM}{R^2}$（與物體質量無關）。高度 $h$ 處 $g_h=\dfrac{GM}{(R+h)^2}$。恆等式 $GM=gR^2$。質量到處不變，重量隨 $g$ 改變。
\item \textbf{圓軌道衛星}：萬有引力＝向心力 $\dfrac{GMm}{r^2}=\dfrac{mv^2}{r}$。軌道速率 $v=\sqrt{\dfrac{GM}{r}}$（外圈較慢，與衛星質量無關）。第一宇宙速度 $v_1=\sqrt{gR}\approx7.9\ \text{km/s}$。同步衛星 $T=24\ \text{h}$，$r\approx4.2\times10^7\ \text{m}$。
\item \textbf{失重}：不是沒有重力，而是人與太空站\textbf{一起自由下落}，支撐力為零。
\item \textbf{克卜勒第三定律（圓軌道推導）}：$T^2=\dfrac{4\pi^2}{GM}r^3\Rightarrow T^2\propto r^3$。比例常數只由中央天體 $M$ 決定；橢圓時 $r$ 換成半長軸 $a$。
\end{itemize}
\end{補充框}

\end{document}
